\chapter{Boolean Analysis — Normal Form Conversion}
%%% generated by the blueprint pipeline from TCSlib.BooleanAnalysis.LMN.NormalFormConversion
\section{Overview}

This module bridges the depth-2 Boolean circuit representation
(\texttt{NAndCircuit} / \texttt{NOrCircuit} built from literals \texttt{BoolCircuit.Lit})
and the DNF/CNF normal forms used by the switching lemma. It defines the
conversion of circuit literals and clauses into switching-lemma terms, and proves
that these conversions preserve evaluation, term width, and the structural
hypotheses (no repeated indices, variable-injectivity) required downstream.

%%%%%%%%%%%%%%%%%%%%%%%%%%%%%%%%%%%%%%%%%%%%%%%%%%%%%%%%%%%%%%%%%%%%%%%%%%%%%%%
\section{Declarations}

\begin{definition}[Convert a circuit literal to a switching-lemma literal]
\lean{BoolCircuit.Lit.toLiteral}
\label{BoolCircuit.Lit.toLiteral}
\leanok
\uses{BoolCircuit.Lit, Literal}
Maps a Boolean-circuit literal $l : \texttt{BoolCircuit.Lit}\,n$ to the corresponding
switching-lemma $\texttt{Literal}\,n$, recording the same variable index and polarity.
\end{definition}

\begin{theorem}[Conversion preserves literal evaluation]
\lean{BoolCircuit.Lit.eval_eq_toLiteral_eval}
\label{BoolCircuit.Lit.eval_eq_toLiteral_eval}
\leanok
\uses{BoolCircuit.Lit, BoolCircuit.Lit.eval, BoolCircuit.Lit.toLiteral, Literal.eval}
\difficulty{2}
Let $l$ be a circuit literal on $n$ variables and let $x : \mathrm{Fin}\,n \to
\mathrm{Bool}$ be an assignment of Boolean values to those variables. Then evaluating
$l$ as a circuit literal at $x$ yields the same Boolean value as first converting $l$ to
the corresponding switching-lemma literal and evaluating that at $x$.
\end{theorem}

\begin{definition}[AND-clause to a term]
\lean{BoolCircuit.NAndCircuit.clauseToTerm}
\statementsource{boolean-ch04-dnf-switching-lmn}{
  pdf-fc447f326910-p001-b004,
  pdf-fc447f326910-p011-b004
}
\label{BoolCircuit.NAndCircuit.clauseToTerm}
\leanok
\uses{BoolCircuit.Lit.toLiteral, BoolCircuit.NAndCircuit, Term}
Sends an AND-clause $\texttt{NAndCircuit}.\mathrm{clause}\,lits$ to the
$\texttt{Term}$ (a conjunction of literals) obtained by mapping each circuit literal
through \texttt{Lit.toLiteral}.
\end{definition}

\begin{definition}[OR-clause to a term]
\lean{BoolCircuit.NOrCircuit.clauseToTerm}
\statementsource{boolean-ch04-dnf-switching-lmn}{
  pdf-fc447f326910-p002-b004,
  pdf-fc447f326910-p011-b004
}
\label{BoolCircuit.NOrCircuit.clauseToTerm}
\leanok
\uses{BoolCircuit.Lit.toLiteral, Term}
Sends an OR-clause $\texttt{NOrCircuit}.\mathrm{clause}\,lits$ to the $\texttt{Term}$
obtained by mapping each circuit literal through \texttt{Lit.toLiteral}; this term plays
the role of one CNF clause (a disjunction of literals).
\end{definition}

\begin{definition}[Depth-2 OR-circuit to a DNF]
\lean{BoolCircuit.NOrCircuit.toDNF}
\statementsource{boolean-ch04-dnf-switching-lmn}{
  pdf-fc447f326910-p011-b004,
  pdf-fc447f326910-p001-b004
}
\label{BoolCircuit.NOrCircuit.toDNF}
\leanok
\uses{BoolCircuit.NAndCircuit.clauseToTerm, DNF}
Converts a depth-2 $\texttt{NOrCircuit}$ (an OR of AND-clauses) into a $\texttt{DNF}$, the
list of terms obtained by converting each AND-clause via \texttt{NAndCircuit.clauseToTerm};
a bare clause maps to the empty DNF.
\end{definition}

\begin{definition}[Depth-2 AND-circuit to a CNF]
\lean{BoolCircuit.NAndCircuit.toCNF}
\statementsource{boolean-ch04-dnf-switching-lmn}{
  pdf-fc447f326910-p011-b004,
  pdf-fc447f326910-p002-b004
}
\label{BoolCircuit.NAndCircuit.toCNF}
\leanok
\uses{BoolCircuit.NAndCircuit, BoolCircuit.NOrCircuit.clauseToTerm, CNF}
Converts a depth-2 $\texttt{NAndCircuit}$ (an AND of OR-clauses) into a $\texttt{CNF}$, the
list of clauses obtained by converting each OR-clause via \texttt{NOrCircuit.clauseToTerm};
a bare clause maps to the empty CNF.
\end{definition}

\begin{theorem}[Conjunction fold of literals equals term evaluation]
\lean{BoolCircuit.foldr_and_lits_eq_term_eval}
\label{BoolCircuit.foldr_and_lits_eq_term_eval}
\leanok
\uses{BoolCircuit.Lit, BoolCircuit.Lit.eval, BoolCircuit.Lit.toLiteral, Term.eval}
\difficulty{3}
Fix a number of variables $n$ and a Boolean assignment $x : \mathrm{Fin}\,n \to
\mathrm{Bool}$. For any finite list of circuit literals $\ell_1, \ldots, \ell_k$ on $n$
variables, the right fold of logical conjunction over their evaluations, taken with
initial value $\mathrm{true}$ — that is, $\ell_1.\mathrm{eval}(x) \wedge
\bigl(\ell_2.\mathrm{eval}(x) \wedge \cdots \wedge (\ell_k.\mathrm{eval}(x) \wedge
\mathrm{true})\bigr)$ — equals the term evaluation, on $x$, of the list obtained by
converting each circuit literal to its corresponding switching-lemma literal.
\end{theorem}

\begin{theorem}[Fold of ORs over literals equals clause evaluation]
\lean{BoolCircuit.foldr_or_lits_eq_clause_eval}
\label{BoolCircuit.foldr_or_lits_eq_clause_eval}
\leanok
\uses{BoolCircuit.Lit, BoolCircuit.Lit.eval, BoolCircuit.Lit.toLiteral, CNF.evalClause}
\difficulty{3}
Let $\ell_1, \dots, \ell_k$ be a list of circuit literals on $n$ variables and let $x :
\mathrm{Fin}\,n \to \mathrm{Bool}$ be an assignment. Then the right fold of disjunction
over the literal evaluations, starting from $\mathrm{false}$,
\[
\ell_1.\mathrm{eval}\,x \;\lor\; \bigl(\ell_2.\mathrm{eval}\,x \;\lor\; \cdots \;\lor\;
(\ell_k.\mathrm{eval}\,x \;\lor\; \mathrm{false})\bigr),
\]
equals the clause evaluation on $x$ of the term obtained by converting each circuit
literal to its switching-lemma literal; that is, both equal $\mathrm{true}$ precisely
when some literal in the list holds under $x$.
\end{theorem}

\begin{theorem}[Depth-2 OR-circuit agrees with its DNF translation]
\lean{BoolCircuit.NOrCircuit.node_eval_eq_toDNF_eval}
\label{BoolCircuit.NOrCircuit.node_eval_eq_toDNF_eval}
\leanok
\uses{BoolCircuit.NAndCircuit, BoolCircuit.NAndCircuit.clauseToTerm, BoolCircuit.NAndCircuit.eval, BoolCircuit.NOrCircuit.eval, BoolCircuit.NOrCircuit.toDNF, BoolCircuit.foldr_and_lits_eq_term_eval, DNF.eval}
\difficulty{4}
Let $cs$ be a list of normal-form AND circuits on $n$ variables, and suppose each member
of $cs$ is a base AND-clause (a conjunction of literals with distinct variable indices)
rather than an AND-node built from OR-subcircuits. Form the normal-form OR circuit whose
children are exactly the circuits in $cs$, so that it computes their disjunction. Then
for every assignment $x : \mathrm{Fin}\,n \to \mathrm{Bool}$, this OR circuit evaluates
at $x$ to the same Boolean value as the DNF formula obtained from it by converting each
child AND-clause into a term.
\end{theorem}

\begin{theorem}[Evaluation of a depth-2 AND circuit matches its CNF]
\lean{BoolCircuit.NAndCircuit.node_eval_eq_toCNF_eval}
\label{BoolCircuit.NAndCircuit.node_eval_eq_toCNF_eval}
\leanok
\uses{BoolCircuit.NAndCircuit.eval, BoolCircuit.NAndCircuit.toCNF, BoolCircuit.NOrCircuit.clauseToTerm, BoolCircuit.NOrCircuit.eval, BoolCircuit.foldr_or_lits_eq_clause_eval, CNF.eval, CNF.evalClause, Term.eval}
\difficulty{4}
Let $cs$ be a list of normal-form OR circuits on $n$ variables, and suppose that every
entry of $cs$ is an OR-clause — that is, a base clause given by a list of literals,
rather than an internal OR node. Then for every assignment $x : \mathrm{Fin}\,n \to
\bbf_2$, evaluating the normal-form AND circuit that conjoins the circuits in $cs$
agrees with evaluating the CNF formula built from it: the value at $x$ of the AND-node
over $cs$ equals the value at $x$ of the CNF whose clauses are obtained by converting
each OR-clause of $cs$ into a disjunction of literals.
\end{theorem}

\begin{theorem}[AND-clause conversion preserves distinctness]
\lean{BoolCircuit.NAndCircuit.clauseToTerm_nodup}
\label{BoolCircuit.NAndCircuit.clauseToTerm_nodup}
\leanok
\uses{BoolCircuit.Lit, BoolCircuit.NAndCircuit.clauseToTerm}
\difficulty{4}
Let $lits$ be a list of Boolean-circuit literals over $n$ variables, and suppose the
list of variable indices $\mathrm{idx}(l)$ taken over the literals $l$ in $lits$ has no
repeated entry. Then the term obtained from the normal-form AND-clause on $lits$ by
mapping each circuit literal to its switching-lemma literal, viewed as a list of
literals, has no repeated entry.
\end{theorem}

\begin{theorem}[Variable-injectivity of a converted AND-clause]
\lean{BoolCircuit.NAndCircuit.clauseToTerm_var_inj}
\label{BoolCircuit.NAndCircuit.clauseToTerm_var_inj}
\leanok
\uses{BoolCircuit.Lit, BoolCircuit.Lit.toLiteral, BoolCircuit.NAndCircuit.clauseToTerm}
\difficulty{4}
Let $lits$ be a list of circuit literals on $n$ variables whose variable indices are
pairwise distinct, and let $C$ be the AND-clause built from $lits$. If $\ell_1$ and
$\ell_2$ are any two literals occurring in the term obtained by converting $C$, and
$\ell_1$ and $\ell_2$ have the same variable, then $\ell_1 = \ell_2$.
\end{theorem}

\begin{theorem}[OR-clause conversion preserves distinctness]
\lean{BoolCircuit.NOrCircuit.clauseToTerm_nodup}
\label{BoolCircuit.NOrCircuit.clauseToTerm_nodup}
\leanok
\uses{BoolCircuit.Lit, BoolCircuit.NAndCircuit.clauseToTerm_nodup, BoolCircuit.NOrCircuit.clauseToTerm}
\difficulty{1}
Let $lits$ be a list of circuit literals in $n$ variables whose underlying variable
indices are pairwise distinct. Then the term obtained by converting the OR-clause on
$lits$ — the list of switching-lemma literals produced by sending each circuit literal
to its counterpart — has no repeated entries.
\end{theorem}

\begin{theorem}[Variable-injectivity of a converted OR-clause]
\lean{BoolCircuit.NOrCircuit.clauseToTerm_var_inj}
\label{BoolCircuit.NOrCircuit.clauseToTerm_var_inj}
\leanok
\uses{BoolCircuit.Lit, BoolCircuit.NAndCircuit.clauseToTerm_var_inj, BoolCircuit.NOrCircuit.clauseToTerm}
\difficulty{1}
Let $n \ge 0$ and let $lits$ be a list of circuit literals over $n$ variables whose
variable indices are pairwise distinct, and form the term obtained by converting each
literal of $lits$ to a switching-lemma literal (recording the same variable and
polarity). Then any two literals of this term that share the same variable $\mathtt{var}
\in \mathrm{Fin}\,n$ are in fact equal.
\end{theorem}

\begin{theorem}[Width of a converted AND-clause]
\lean{BoolCircuit.NAndCircuit.clauseToTerm_width}
\label{BoolCircuit.NAndCircuit.clauseToTerm_width}
\leanok
\uses{BoolCircuit.Lit, BoolCircuit.NAndCircuit.clauseToTerm, Term.width}
\difficulty{1}
Fix $n$, and let $lits$ be a list of circuit literals on $n$ variables whose variable
indices are pairwise distinct. Form the normal-form AND-clause on this list, and convert
it into a term by replacing each circuit literal with the corresponding switching-lemma
literal. Then the width of that term — its number of literals — equals the length of
$lits$.
\end{theorem}

\begin{theorem}[Width of the term from an OR-clause]
\lean{BoolCircuit.NOrCircuit.clauseToTerm_width}
\label{BoolCircuit.NOrCircuit.clauseToTerm_width}
\leanok
\uses{BoolCircuit.Lit, BoolCircuit.NOrCircuit.clauseToTerm, Term.width}
\difficulty{1}
Let $\mathrm{lits}$ be a list of circuit literals over $n$ variables whose variable
indices are pairwise distinct. Then the width of the term obtained from the OR-clause on
$\mathrm{lits}$ — that is, the term formed by converting each circuit literal to its
corresponding literal — equals the number of literals in $\mathrm{lits}$.
\end{theorem}

\begin{theorem}[Width bound for the DNF of a disjunction of AND-clauses]
\lean{BoolCircuit.NOrCircuit.toDNF_width_bounded}
\label{BoolCircuit.NOrCircuit.toDNF_width_bounded}
\leanok
\uses{BoolCircuit.Lit, BoolCircuit.NAndCircuit, BoolCircuit.NAndCircuit.clauseToTerm_width, BoolCircuit.NOrCircuit.toDNF, DNF.width}
\difficulty{4}
Let $cs$ be a finite list of normal-form AND circuits on $n$ variables, and let $w \in
\bbn$. Suppose every circuit in $cs$ is a base clause — a list of literals whose
variable indices are pairwise distinct — and that each such clause contains at most $w$
literals. Then the DNF obtained from the normal-form OR circuit whose disjuncts are the
circuits in $cs$, by converting each AND-clause into a term, has width at most $w$; that
is, every term in it has at most $w$ literals.
\end{theorem}

\begin{theorem}[Width bound for the CNF of an AND of OR-clauses]
\lean{BoolCircuit.NAndCircuit.toCNF_width_bounded}
\label{BoolCircuit.NAndCircuit.toCNF_width_bounded}
\leanok
\uses{BoolCircuit.Lit, BoolCircuit.NAndCircuit.toCNF, BoolCircuit.NOrCircuit.clauseToTerm_width, CNF.width}
\difficulty{4}
Fix $n \in \bbn$, let $cs$ be a list of normal-form OR circuits on $n$ variables, and
let $w \in \bbn$. Suppose every circuit in $cs$ is an OR-clause on some list of literals
whose variable indices are pairwise distinct and whose length is at most $w$. Form the
normal-form AND circuit that conjoins all the members of $cs$, and convert that circuit
to a CNF formula by turning each OR-clause into a clause (a disjunction of its
literals). Then this CNF formula has width at most $w$; equivalently, each of its
clauses contains at most $w$ literals.
\end{theorem}

\begin{theorem}[DNF terms of a normal-form OR circuit are duplicate-free]
\lean{BoolCircuit.NOrCircuit.toDNF_terms_nodup}
\statementsource{boolean-ch04-dnf-switching-lmn}{
  pdf-fc447f326910-p001-b004,
  pdf-fc447f326910-p012-b001
}
\label{BoolCircuit.NOrCircuit.toDNF_terms_nodup}
\leanok
\uses{BoolCircuit.Lit, BoolCircuit.NAndCircuit, BoolCircuit.NAndCircuit.clauseToTerm_nodup, BoolCircuit.NOrCircuit.toDNF}
\difficulty{3}
Let $cs$ be a list of normal-form AND circuits on $n$ variables, and suppose every
element of $cs$ is a clause — that is, each is given by a list of circuit literals whose
variable indices are pairwise distinct. Consider the normal-form OR circuit that
combines the members of $cs$, and let its associated DNF formula be the list of terms
obtained by converting each clause into a term (mapping every circuit literal to the
corresponding switching-lemma literal with the same variable and polarity). Then every
term of this DNF is duplicate-free: no term contains the same literal twice.
\end{theorem}

\begin{theorem}[Variable-injectivity of literals in the DNF of an OR-of-clauses]
\lean{BoolCircuit.NOrCircuit.toDNF_var_inj}
\label{BoolCircuit.NOrCircuit.toDNF_var_inj}
\leanok
\uses{BoolCircuit.Lit, BoolCircuit.NAndCircuit, BoolCircuit.NAndCircuit.clauseToTerm, BoolCircuit.NAndCircuit.clauseToTerm_var_inj, BoolCircuit.NOrCircuit.toDNF}
\difficulty{3}
Fix $n \in \bbn$ and a list $c_1, \dots, c_m$ of normal-form AND circuits on $n$
variables, and suppose each $c_i$ is a base AND-clause — that is, $c_i =
\mathrm{clause}(\ell_1, \dots, \ell_{k_i})$ for some list of circuit literals whose
variable indices are pairwise distinct. Let $F$ be the DNF formula obtained from the OR
circuit $\mathrm{node}(c_1, \dots, c_m)$ by converting each clause $c_i$ into the term
formed by its literals. Then in every term of $F$, any two literals that share the same
variable are identical.
\end{theorem}

\begin{theorem}[Duplicate-free clauses in the CNF of an AND-node]
\lean{BoolCircuit.NAndCircuit.toCNF_terms_nodup}
\label{BoolCircuit.NAndCircuit.toCNF_terms_nodup}
\leanok
\uses{BoolCircuit.Lit, BoolCircuit.NAndCircuit.toCNF, BoolCircuit.NOrCircuit.clauseToTerm_nodup}
\difficulty{3}
Let $cs$ be a list of normal-form OR circuits on $n$ variables, and suppose every member
of $cs$ is an OR-clause $\mathrm{clause}\,(lits)$ whose literals have pairwise distinct
variable indices. Form the depth-2 AND circuit that conjoins the circuits in $cs$, and
pass to the CNF formula obtained by turning each of these OR-clauses into a clause (a
disjunction of literals). Then every clause of this CNF is a list of literals with no
repeated entry.
\end{theorem}

\begin{theorem}[Variable-injectivity within the clauses of an AND-node's CNF]
\lean{BoolCircuit.NAndCircuit.toCNF_var_inj}
\label{BoolCircuit.NAndCircuit.toCNF_var_inj}
\leanok
\uses{BoolCircuit.Lit, BoolCircuit.NAndCircuit.toCNF, BoolCircuit.NOrCircuit.clauseToTerm, BoolCircuit.NOrCircuit.clauseToTerm_var_inj}
\difficulty{3}
Let $cs$ be a list of normal-form OR circuits on $n$ variables, and suppose every member
of $cs$ is a base OR-clause: for each $c \in cs$ there is a list of circuit literals
$lits$, whose variable indices are pairwise distinct, with $c = \mathrm{clause}\,lits$.
Form the CNF formula of the AND-node over $cs$ by converting each of these OR-clauses
into a term. Then within any single clause of this CNF, any two literals that share the
same variable are identical — in particular they carry the same negation flag.
\end{theorem}
